% Bagian: computability
% Bab: recursive-functions
% Subbagian: normal-form

\documentclass[../../../include/open-logic-section]{subfiles}

\begin{document}

\olfileid[id]{cmp}{rec}{nft}
\olsection{Teorema Bentuk Normal}

\begin{thm}[Teorema Bentuk Normal Kleene]
\ollabel{thm:kleene-nf}
Terdapat suatu relasi rekursif primitif $T(e,x,s)$ dan suatu fungsi rekursif
primitif $U(s)$ dengan sifat berikut: jika $f$ merupakan sebarang fungsi
rekursif parsial uner, maka bagi suatu~$e$,
\[
f(x) \simeq U(\umin{s}{T(e, x, s)})
\]
untuk setiap~$x$.
\end{thm}

\begin{explain}
Bukti teorema bentuk normal cukup rumit, tetapi gagasan dasarnya sederhana.
Setiap fungsi rekursif parsial mempunyai suatu \emph{indeks}~$e$, yang secara
intuitif merupakan bilangan yang mengodekan program atau definisinya. Jika
$f(x)\fdefined$, komputasi itu dapat dicatat secara sistematis dan dikodekan
oleh suatu bilangan~$s$, dan fakta bahwa $s$ mengodekan komputasi~$f$ pada
masukan~$x$ dapat diperiksa secara rekursif primitif hanya dengan menggunakan
$x$ dan definisi~$e$. Oleh karena itu, relasi~$T$, ``fungsi berindeks~$e$
mempunyai suatu komputasi bagi masukan~$x$, dan $s$ mengodekan komputasi
tersebut,'' bersifat rekursif primitif. Dari catatan lengkap komputasi~$s$,
``hasil akhir''~$s$, yaitu nilai~$f(x)$, juga dapat diperoleh dari~$s$ secara
rekursif primitif.

Teorema bentuk normal menunjukkan bahwa hanya satu pencarian tak berbatas yang
diperlukan dalam definisi sebarang fungsi rekursif parsial. Pada dasarnya,
kita dapat mencari di antara semua bilangan sampai menemukan bilangan yang
mengodekan suatu komputasi dari fungsi berindeks~$e$ bagi masukan~$x$. Kita
dapat menggunakan bilangan~$e$ sebagai ``nama'' fungsi rekursif parsial, dan
menulis $\cfind{e}$ bagi fungsi~$f$ yang didefinisikan oleh persamaan dalam
teorema. Perhatikan bahwa suatu fungsi rekursif parsial dapat mempunyai lebih
dari satu indeks---bahkan, setiap fungsi rekursif parsial mempunyai tak hingga
banyak indeks.
\end{explain}
\end{document}
